\chapter{Taylor Approximation}

The previous chapter closed with a formula for the Taylor polynomial;
this chapter turns it into a working tool. Most functions that appear in
economics (utility, production, growth) are not polynomials, yet
polynomials are the easiest objects to evaluate and to reason about.
Taylor's theorem says that near a chosen base point every sufficiently
smooth function \emph{is} a polynomial, up to a controllable error. This
single idea underlies the linearisation of models around a steady state,
the log-linearisation trick $\ln(1+g)\approx g$ for growth rates, the
second-order conditions of optimisation, and quick numerical estimates
such as $\sqrt{4.1}$ or $\mathrm{e}^{0.1}$ without a calculator.

\section{Core Rules}

\begin{rulebox}{Taylor polynomial of order $n$}
Let $f$ be $n$ times differentiable at $a$. The \emph{Taylor polynomial}
of order $n$ of $f$ at $a$ is
\[
  T_{n}(x)=\sum_{k=0}^{n}\frac{f^{(k)}(a)}{k!}(x-a)^{k}
          =f(a)+f'(a)(x-a)+\frac{f''(a)}{2!}(x-a)^{2}+\cdots
           +\frac{f^{(n)}(a)}{n!}(x-a)^{n}.
\]
It is the \emph{unique} polynomial of degree $\leq n$ whose value and
first $n$ derivatives at $a$ all agree with those of $f$:
$T_{n}^{(k)}(a)=f^{(k)}(a)$ for $k=0,1,\ldots,n$.
\end{rulebox}

\begin{rulebox}{Linear and quadratic approximation ($n=1,2$)}
\[
  T_{1}(x)=f(a)+f'(a)(x-a),
  \qquad
  T_{2}(x)=f(a)+f'(a)(x-a)+\tfrac12 f''(a)(x-a)^{2}.
\]
$T_{1}$ is the tangent line at $a$ (the \emph{linearisation} of $f$);
$T_{2}$ is the best-fitting parabola, matching value, slope and
curvature at $a$. Each extra order typically buys one further correct
decimal place near $a$.
\end{rulebox}

\begin{rulebox}{Lagrange remainder and error bound}
If $f$ is $(n+1)$ times differentiable on an interval containing $a$ and
$x$, then
\[
  f(x)=T_{n}(x)+R_{n}(x),
  \qquad
  R_{n}(x)=\frac{f^{(n+1)}(\xi)}{(n+1)!}(x-a)^{n+1}
\]
for some $\xi$ strictly between $a$ and $x$. Since $\xi$ is unknown,
bound it: if $|f^{(n+1)}(t)|\leq M$ for all $t$ between $a$ and $x$,
then
\[
  \bigl|f(x)-T_{n}(x)\bigr|\leq\frac{M}{(n+1)!}\,|x-a|^{n+1}.
\]
The error shrinks like $|x-a|^{n+1}$: halving the distance to $a$
divides the error of $T_{n}$ by about $2^{n+1}$.
\end{rulebox}

\begin{rulebox}{Maclaurin case ($a=0$): standard expansions}
The Taylor polynomial at $a=0$ is called the \emph{Maclaurin}
polynomial. The classical expansions are
\[
\begin{aligned}
  \mathrm{e}^{x}&=1+x+\frac{x^{2}}{2!}+\frac{x^{3}}{3!}+\cdots
                  +\frac{x^{n}}{n!}+\cdots
                  &&\text{(all }x\text{)},\\
  \ln(1+x)&=x-\frac{x^{2}}{2}+\frac{x^{3}}{3}-\frac{x^{4}}{4}+\cdots
                  &&(-1<x\leq 1),\\
  (1+x)^{k}&=1+kx+\frac{k(k-1)}{2!}x^{2}
                  +\frac{k(k-1)(k-2)}{3!}x^{3}+\cdots
                  &&(|x|<1),\\
  \sin(x)&=x-\frac{x^{3}}{3!}+\frac{x^{5}}{5!}-\cdots
                  &&\text{(all }x\text{)},\\
  \cos(x)&=1-\frac{x^{2}}{2!}+\frac{x^{4}}{4!}-\cdots
                  &&\text{(all }x\text{)}.
\end{aligned}
\]
The first-order case $\ln(1+x)\approx x$ is the workhorse of
log-linearisation in economics: a small growth rate $g$ and the log
change $\ln(1+g)$ are interchangeable to first order.
\end{rulebox}

\begin{rulebox}{Quadratic approximation in economics}
Near a stationary point $a$ (where $f'(a)=0$) the linear term vanishes
and
\[
  f(x)\approx f(a)+\tfrac12 f''(a)(x-a)^{2}:
\]
locally, $f$ \emph{is} a parabola sitting on top of $f(a)$. If
$f''(a)<0$ the parabola opens downward and $a$ is a local maximum; if
$f''(a)>0$ it opens upward and $a$ is a local minimum. This is exactly
where the second-order condition comes from. In several variables the
role of $f''(a)$ is played by the Hessian, and the sign condition
becomes definiteness.
\end{rulebox}

\begin{rulebox}{Caution: validity is local}
A Taylor polynomial is a \emph{local} approximation: the error grows
like $|x-a|^{n+1}$, so accuracy is excellent near $a$ and can be
useless far away. For $\ln(1+x)$ and $(1+x)^{k}$ the expansions are
only valid on $(-1,1]$ resp.\ $(-1,1)$: outside, adding more terms
makes things \emph{worse}. Even inside: $T_{1}(x)=x$ for $\ln(1+x)$
gives $\ln(2)\approx 1$ against the true $0.693$, a $44\%$ error,
versus $0.5\%$ at $x=0.1$. Rule of thumb: expand around a base point
close to where you evaluate, and use the remainder bound to check.
\end{rulebox}

\section{Worked Example}

\begin{worked}{Expand $\ln(1+x)$ to third order about $0$, bound the error at $x=0.5$, and read the economics off the second-order term}
\wstep{Step 1. Collect the derivatives at the expansion point.}
\[
  f(x)=\ln(1+x),\quad
  f'=\frac{1}{1+x},\quad
  f''=\frac{-1}{(1+x)^{2}},\quad
  f'''=\frac{2}{(1+x)^{3}},\quad
  f^{(4)}=\frac{-6}{(1+x)^{4}} .
\]
At $a=0$: $f(0)=0$, $f'(0)=1$, $f''(0)=-1$, $f'''(0)=2$, and
$f^{(4)}(0)=-6$.
\wtrap{Evaluate \emph{after} differentiating, never before. Substituting
$x=0$ into $f'$ and then differentiating the resulting constant gives zeros
all the way down --- a surprisingly frequent way to produce the polynomial
$P_{3}(x)=x$.}

\wstep{Step 2. Assemble the Taylor polynomial.}
\[
  P_{n}(x)=\sum_{k=0}^{n}\frac{f^{(k)}(a)}{k!}(x-a)^{k}
  \quad\Longrightarrow\quad
  P_{3}(x)=0+x-\frac{1}{2!}x^{2}+\frac{2}{3!}x^{3}
  =\boxed{x-\frac{x^{2}}{2}+\frac{x^{3}}{3}}
\]
The factorials are what turn the raw derivatives $1,-1,2$ into the
coefficients $1,-\tfrac12,\tfrac13$.
\wtrap{Do not omit $k!$, and do not omit $(x-a)^{k}$ when $a\neq0$. Expanding
about $a=2$ and writing powers of $x$ rather than of $x-2$ is the standard
error, and it silently produces a polynomial that is wrong everywhere.}

\wstep{Step 3. Evaluate at $x=0.5$ and compare.}
\[
  P_{3}(0.5)=0.5-0.125+0.041\overline{6}=0.4166\overline{6},
  \qquad
  \ln(1.5)=0.405465 ,
\]
so the actual error is $0.011202$.

\wstep{Step 4. Bound the error \emph{before} knowing the true value.}
The Lagrange remainder is
\[
  R_{3}(x)=\frac{f^{(4)}(\xi)}{4!}x^{4}
  \qquad\text{for some }\xi\text{ between }0\text{ and }x .
\]
$\xi$ is unknown, so take the worst case. On $[0,0.5]$,
\[
  \bigl|f^{(4)}(\xi)\bigr|=\frac{6}{(1+\xi)^{4}}\le\frac{6}{1^{4}}=6
  \qquad(\text{largest at }\xi=0),
\]
\[
  |R_{3}(0.5)|\le\frac{6}{24}(0.5)^{4}=0.25\cdot0.0625=0.015625 .
\]
And indeed $0.011201\le0.015625$. \checkmark
\wtrap{Maximise $|f^{(n+1)}|$ over the whole interval between $a$ and $x$ ---
do not substitute $\xi=x$ or $\xi=a$ and call it a bound. Here $1/(1+\xi)^{4}$
is decreasing, so the worst case is at the \emph{left} end; for an
increasing derivative it would be at the right. Check which end before
committing.}

\wstep{Step 5. Note how fast the accuracy degrades with distance.}
At $x=0.1$ the same bound gives $\tfrac14(0.1)^{4}=0.000025$; at $x=0.5$ it
is $625$ times larger. The error carries the factor $x^{4}$, so halving the
distance to the expansion point cuts the error sixteenfold. At $x=1$ the
series still converges, but at $x=1.1$ it diverges outright --- the
expansion of $\ln(1+x)$ is valid only for $|x|<1$.
\wtrap{A Taylor polynomial is a \textbf{local} statement. It is excellent
near $a$, useless far from it, and there is no warning in the formula --- 
$P_{3}(5)=5-12.5+41.7=34.2$, while $\ln6=1.79$. Always ask how far from $a$
you are being asked to work.}

\wstep{Step 6. The first two terms are the economics.}
Keeping $n=1$ gives the \textbf{linear approximation} $\ln(1+x)\approx x$:
the familiar ``a log difference is a percentage change'', accurate to about
half a percentage point for $|x|\le0.1$. Keeping $n=2$ gives
\[
  f(x)\approx f(a)+f'(a)(x-a)+\tfrac12 f''(a)(x-a)^{2},
\]
whose quadratic term carries the \textbf{curvature}. In expected-utility
theory that term is exactly what produces the risk premium: a concave $u$
has $u''<0$, so the quadratic term is negative, and its size relative to
$u'$ is the Arrow--Pratt coefficient $-u''/u'$. Risk aversion \emph{is} the
second-order term of a Taylor expansion.

\wstep{Step 7. The expansions worth knowing by heart.}
\[
  \mathrm{e}^{x}=1+x+\tfrac{x^{2}}{2}+\tfrac{x^{3}}{6}+\cdots,
  \qquad
  \ln(1+x)=x-\tfrac{x^{2}}{2}+\tfrac{x^{3}}{3}-\cdots,
\]
\[
  (1+x)^{\alpha}=1+\alpha x+\tfrac{\alpha(\alpha-1)}{2}x^{2}+\cdots,
  \qquad
  \frac{1}{1-x}=1+x+x^{2}+\cdots
\]
The last is the geometric series of Volume~4, which is why a perpetuity
formula and a Taylor expansion are the same computation twice.
\end{worked}

\section{Practice Problems}

\begin{problems}
  \item Find $T_{1}$ of $f(x)=\ln(x)$ at $a=1$.
  \item Find $T_{1}$ of $f(x)=\sqrt{x}$ at $a=4$.
  \item Find $T_{1}$ of $f(x)=\dfrac{1}{x}$ at $a=1$.
  \item Find $T_{1}$ of $f(x)=x^{3}-2x$ at $a=2$.
  \item Find $T_{1}$ of $f(x)=\mathrm{e}^{2x}$ at $a=\ln 2$.
  \item Find $T_{2}$ of $f(x)=\mathrm{e}^{x}$ at $a=0$ and use it to
        approximate $\mathrm{e}^{0.1}$. Compare with the true value
        $\mathrm{e}^{0.1}=1.105171$.
  \item Find $T_{2}$ of $f(x)=\ln(x)$ at $a=1$ and use it to
        approximate $\ln(1.05)$. Compare with the true value
        $\ln(1.05)=0.048790$.
  \item Find $T_{2}$ of $f(x)=\sqrt{x}$ at $a=4$ and use it to
        approximate $\sqrt{4.1}$. Compare with the true value
        $\sqrt{4.1}=2.024846$.
  \item Find $T_{2}$ of $f(x)=\dfrac{1}{1-x}$ at $a=0$.
  \item Find $T_{2}$ of $f(x)=x\ln(x)$ at $a=1$.
  \item Find $T_{2}$ of $f(x)=\mathrm{e}^{-x^{2}}$ at $a=0$.
  \item Find $T_{3}$ of $f(x)=\ln(1+x)$ at $a=0$ and use it to
        approximate $\ln(0.9)$. Compare with the true value
        $\ln(0.9)=-0.105361$.
  \item Use the binomial expansion to find $T_{2}$ of
        $f(x)=(1+x)^{1/2}$ at $a=0$, and use it to approximate
        $\sqrt{1.1}$. Compare with the true value
        $\sqrt{1.1}=1.048809$.
  \item Find the Maclaurin polynomial $T_{4}$ of $f(x)=\cos(x)$ and use
        it to approximate $\cos(0.2)$. Compare with the true value
        $\cos(0.2)=0.980067$.
  \item A firm produces $Q(L)=10\sqrt{L}$ units with $L$ hours of
        labour. Find $T_{2}$ of $Q$ at $a=100$ and use it to
        approximate output at $L=104$. Compare with the true value
        $Q(104)=101.980390$.
  \item (Log-linearisation.) Output grows from $Y_{0}$ to
        $Y_{1}=Y_{0}(1+g)$. Show, using $T_{1}$ of $\ln(1+x)$ at $0$,
        that $\ln Y_{1}-\ln Y_{0}\approx g$, and quantify the
        approximation for $g=0.03$ given $\ln(1.03)=0.029559$.
  \item Use the Lagrange remainder to bound the error of the $T_{1}$
        approximation of $\sqrt{4.1}$ at $a=4$ (Problem 2), and compare
        the bound with the actual error.
  \item Use the Lagrange remainder to bound the error of the $T_{2}$
        approximation of $\mathrm{e}^{0.1}$ at $a=0$ (Problem 6), and
        compare the bound with the actual error.
  \item Find the smallest $n$ such that the Maclaurin polynomial
        $T_{n}$ of $\mathrm{e}^{x}$ approximates $\mathrm{e}^{0.5}$
        with error at most $10^{-4}$. (Use $\mathrm{e}^{0.5}<2$ in the
        remainder bound.)
  \item Prove, using $T_{1}$ with Lagrange remainder: if $f'(a)=0$ and
        $f''$ is continuous with $f''(a)<0$, then $a$ is a strict local
        maximum of $f$.
\end{problems}

\section{Solutions}

\begin{solutions}
  \item $f(1)=\ln(1)=0$; $f'(x)=1/x$, so $f'(1)=1$. Hence
        \[
          T_{1}(x)=0+1\cdot(x-1)=x-1 .
        \]

  \item $f(4)=2$; $f'(x)=\dfrac{1}{2\sqrt{x}}$, so
        $f'(4)=\dfrac{1}{2\cdot 2}=\dfrac14$. Hence
        \[
          T_{1}(x)=2+\tfrac14(x-4) .
        \]

  \item $f(1)=1$; $f'(x)=-1/x^{2}$, so $f'(1)=-1$. Hence
        \[
          T_{1}(x)=1-(x-1)=2-x .
        \]

  \item $f(2)=8-4=4$; $f'(x)=3x^{2}-2$, so $f'(2)=12-2=10$. Hence
        \[
          T_{1}(x)=4+10(x-2) .
        \]

  \item $f(\ln 2)=\mathrm{e}^{2\ln 2}=(\mathrm{e}^{\ln 2})^{2}=4$;
        $f'(x)=2\mathrm{e}^{2x}$, so $f'(\ln 2)=2\cdot 4=8$. Hence
        \[
          T_{1}(x)=4+8(x-\ln 2) .
        \]

  \item $f^{(k)}(x)=\mathrm{e}^{x}$ for every $k$, so
        $f(0)=f'(0)=f''(0)=1$ and
        \[
          T_{2}(x)=1+x+\frac{x^{2}}{2} .
        \]
        At $x=0.1$: $T_{2}(0.1)=1+0.1+\tfrac{0.01}{2}=1.105$. True
        value $1.105171$; the error is about $1.7\cdot 10^{-4}$.

  \item $f(1)=0$; $f'(x)=1/x$, $f'(1)=1$; $f''(x)=-1/x^{2}$,
        $f''(1)=-1$. Hence
        \[
          T_{2}(x)=(x-1)-\tfrac12(x-1)^{2} .
        \]
        At $x=1.05$:
        $T_{2}(1.05)=0.05-\tfrac12(0.0025)=0.05-0.00125=0.04875$. True
        value $0.048790$; the error is about $4.0\cdot 10^{-5}$.

  \item From Problem 2, $f(4)=2$ and $f'(4)=\tfrac14$. Further
        $f''(x)=-\tfrac14 x^{-3/2}$, so
        $f''(4)=-\tfrac14\cdot\tfrac18=-\tfrac{1}{32}$ and the
        quadratic coefficient is
        $\tfrac12 f''(4)=-\tfrac{1}{64}$:
        \[
          T_{2}(x)=2+\tfrac14(x-4)-\tfrac{1}{64}(x-4)^{2} .
        \]
        At $x=4.1$:
        \[
          T_{2}(4.1)=2+0.025-\tfrac{0.01}{64}
                    =2.025-0.00015625=2.02484375 .
        \]
        True value $2.024846$; the error is about $1.9\cdot 10^{-6}$
        (the linear estimate $2.025$ alone is off by
        $1.5\cdot 10^{-4}$).

  \item $f(0)=1$; $f'(x)=\dfrac{1}{(1-x)^{2}}$, $f'(0)=1$;
        $f''(x)=\dfrac{2}{(1-x)^{3}}$, $f''(0)=2$, so
        $\tfrac12 f''(0)=1$. Hence
        \[
          T_{2}(x)=1+x+x^{2} ,
        \]
        the start of the geometric series $1+x+x^{2}+\cdots$.

  \item $f(1)=1\cdot\ln(1)=0$; $f'(x)=\ln(x)+1$, $f'(1)=1$;
        $f''(x)=1/x$, $f''(1)=1$. Hence
        \[
          T_{2}(x)=(x-1)+\tfrac12(x-1)^{2} .
        \]

  \item $f(0)=1$; $f'(x)=-2x\mathrm{e}^{-x^{2}}$, $f'(0)=0$;
        $f''(x)=(-2+4x^{2})\mathrm{e}^{-x^{2}}$, $f''(0)=-2$. Hence
        \[
          T_{2}(x)=1-x^{2} .
        \]
        (Equivalently: substitute $u=-x^{2}$ into
        $\mathrm{e}^{u}\approx 1+u$.)

  \item $f(0)=0$; $f'(x)=\dfrac{1}{1+x}$, $f'(0)=1$;
        $f''(x)=-\dfrac{1}{(1+x)^{2}}$, $f''(0)=-1$;
        $f'''(x)=\dfrac{2}{(1+x)^{3}}$, $f'''(0)=2$. The coefficients
        are $1$, $-\tfrac{1}{2!}=-\tfrac12$ and
        $\tfrac{2}{3!}=\tfrac13$:
        \[
          T_{3}(x)=x-\frac{x^{2}}{2}+\frac{x^{3}}{3} .
        \]
        At $x=-0.1$:
        \[
          T_{3}(-0.1)=-0.1-\frac{0.01}{2}+\frac{-0.001}{3}
                     =-0.1-0.005-0.000333=-0.105333 .
        \]
        True value $-0.105361$; the error is about $2.7\cdot 10^{-5}$.

  \item With $k=\tfrac12$ in the binomial expansion: the linear
        coefficient is $k=\tfrac12$ and the quadratic coefficient is
        \[
          \frac{k(k-1)}{2!}
            =\frac{\tfrac12\cdot\bigl(-\tfrac12\bigr)}{2}
            =-\frac18 ,
        \]
        so
        \[
          T_{2}(x)=1+\frac{x}{2}-\frac{x^{2}}{8} .
        \]
        At $x=0.1$: $T_{2}(0.1)=1+0.05-0.00125=1.04875$. True value
        $1.048809$; the error is about $5.9\cdot 10^{-5}$.

  \item Derivatives of $\cos$ at $0$ cycle through
        $1,\,0,\,-1,\,0,\,1$ for $k=0,\ldots,4$, so
        \[
          T_{4}(x)=1-\frac{x^{2}}{2!}+\frac{x^{4}}{4!}
                  =1-\frac{x^{2}}{2}+\frac{x^{4}}{24} .
        \]
        At $x=0.2$: $(0.2)^{2}=0.04$ and $(0.2)^{4}=0.0016$, so
        \[
          T_{4}(0.2)=1-0.02+\frac{0.0016}{24}
                    =1-0.02+0.000067=0.980067 ,
        \]
        agreeing with the true value $0.980067$ to all six decimals
        shown (the next term is of size $\tfrac{(0.2)^{6}}{720}\approx
        9\cdot 10^{-8}$).

  \item $Q(100)=10\cdot 10=100$; $Q'(L)=5L^{-1/2}$,
        $Q'(100)=\tfrac{5}{10}=0.5$; $Q''(L)=-\tfrac52 L^{-3/2}$,
        $Q''(100)=-\tfrac{2.5}{1000}=-0.0025$, so
        $\tfrac12 Q''(100)=-0.00125$. Hence
        \[
          T_{2}(L)=100+0.5\,(L-100)-0.00125\,(L-100)^{2} .
        \]
        At $L=104$ (so $L-100=4$):
        \[
          T_{2}(104)=100+0.5\cdot 4-0.00125\cdot 16=102-0.02=101.98 .
        \]
        True value $101.980390$; the error is about
        $3.9\cdot 10^{-4}$, four extra labour hours raise output by
        just under $2$ units, the shortfall against the linear estimate
        $102$ reflecting diminishing returns ($Q''<0$).

  \item Taking logs, $\ln Y_{1}=\ln Y_{0}+\ln(1+g)$. For
        $f(x)=\ln(1+x)$ at $a=0$: $f(0)=0$ and $f'(0)=1$, so
        $T_{1}(x)=x$ and
        \[
          \ln Y_{1}-\ln Y_{0}=\ln(1+g)\approx g .
        \]
        For $g=0.03$ the log change is $\ln(1.03)=0.029559$, while the
        growth rate is $0.03$: the approximation is off by about
        $4.4\cdot 10^{-4}$, i.e.\ by roughly $1.5\%$ of the quantity
        itself. (The second-order term $-\tfrac12 g^{2}=-0.00045$
        accounts for almost all of the gap.)

  \item Here $n=1$, $a=4$, $x=4.1$ and $f''(t)=-\tfrac14 t^{-3/2}$.
        For $t\in[4,4.1]$, $t^{-3/2}$ is decreasing, so
        \[
          |f''(t)|\leq\tfrac14\cdot 4^{-3/2}=\tfrac14\cdot\tfrac18
                 =\tfrac{1}{32}=:M .
        \]
        The remainder bound gives
        \[
          |R_{1}(4.1)|\leq\frac{M}{2!}\,(0.1)^{2}
            =\frac{1}{32}\cdot\frac{0.01}{2}
            =\frac{0.01}{64}\approx 1.56\cdot 10^{-4} .
        \]
        Actual error: $T_{1}(4.1)=2.025$ versus $\sqrt{4.1}=2.024846$,
        i.e.\ $1.54\cdot 10^{-4}$, just inside the bound.

  \item Here $n=2$, $a=0$, $x=0.1$ and
        $f'''(t)=\mathrm{e}^{t}$. For $t\in(0,0.1)$,
        $\mathrm{e}^{t}<\mathrm{e}^{0.1}<1.2=:M$, so
        \[
          |R_{2}(0.1)|\leq\frac{M}{3!}\,(0.1)^{3}
            =\frac{1.2\cdot 0.001}{6}=2\cdot 10^{-4} .
        \]
        Actual error: $T_{2}(0.1)=1.105$ versus
        $\mathrm{e}^{0.1}=1.105171$, i.e.\ $1.7\cdot 10^{-4}$,
        consistent with the bound.

  \item With $f(x)=\mathrm{e}^{x}$, $a=0$, $x=0.5$:
        $f^{(n+1)}(\xi)=\mathrm{e}^{\xi}\leq\mathrm{e}^{0.5}<2$ for
        $\xi\in(0,0.5)$, so
        \[
          |R_{n}(0.5)|\leq\frac{2\,(0.5)^{n+1}}{(n+1)!} .
        \]
        Tabulate:
        \[
          n=3:\ \frac{2\cdot 0.0625}{24}\approx 5.2\cdot 10^{-3},
          \qquad
          n=4:\ \frac{2\cdot 0.03125}{120}\approx 5.2\cdot 10^{-4},
        \]
        \[
          n=5:\ \frac{2\cdot 0.015625}{720}\approx 4.3\cdot 10^{-5} .
        \]
        The first value below $10^{-4}$ occurs at $n=5$. Check with the
        true value $\mathrm{e}^{0.5}=1.648721$: $T_{4}(0.5)=1.648438$
        (error $2.8\cdot 10^{-4}$, indeed too large) while
        $T_{5}(0.5)=1.648698$ (error $2.3\cdot 10^{-5}$). So $n=5$.

  \item By Taylor's theorem with $n=1$: for each $x$ near $a$ there is
        $\xi$ between $a$ and $x$ with
        \[
          f(x)=f(a)+f'(a)(x-a)+\tfrac12 f''(\xi)(x-a)^{2}
              =f(a)+\tfrac12 f''(\xi)(x-a)^{2},
        \]
        using $f'(a)=0$. Since $f''$ is continuous and $f''(a)<0$,
        there is a neighbourhood $U$ of $a$ on which $f''<0$. For
        $x\in U$, $x\neq a$, the point $\xi$ also lies in $U$, so
        $f''(\xi)<0$ while $(x-a)^{2}>0$, giving
        \[
          f(x)-f(a)=\tfrac12 f''(\xi)(x-a)^{2}<0 .
        \]
        Hence $f(x)<f(a)$ for all $x\in U\setminus\{a\}$: $a$ is a
        strict local maximum. (This is the second-order condition in
        its cleanest form: near a stationary point the function
        \emph{is} its osculating parabola, up to higher-order terms.)
\end{solutions}

\section{Additional Problems}

\begin{problems}
  \item Find $T_{1}$ of $f(x)=(1+x)^{5}$ at $a=0$.
  \item Find $T_{1}$ of $f(x)=\sin(x)$ at $a=0$ and use it to
        approximate $\sin(0.1)$. Compare with the true value
        $\sin(0.1)=0.099833$.
  \item Find $T_{2}$ of $f(x)=x^{4}$ at $a=1$.
  \item Find $T_{2}$ of $f(x)=\dfrac{1}{x}$ at $a=2$.
  \item Find $T_{3}$ of $f(x)=\mathrm{e}^{x}$ at $a=0$ and use it to
        approximate $\mathrm{e}^{0.2}$. Compare with the true value
        $\mathrm{e}^{0.2}=1.221403$.
  \item Find $T_{2}$ of $f(x)=\ln(1+2x)$ at $a=0$.
  \item (Cobb--Douglas.) For $f(K)=K^{0.3}$, find $T_{2}$ at $a=1$ and
        use it to approximate $1.1^{0.3}$. Compare with the true value
        $1.1^{0.3}=1.029006$.
  \item Use the Lagrange remainder to bound the error of approximating
        $\ln(1.05)$ by $T_{1}(x)=x-1$ of $\ln(x)$ at $a=1$, and compare
        the bound with the actual error (true value
        $\ln(1.05)=0.048790$).
  \item Find the smallest $n$ such that the Maclaurin polynomial
        $T_{n}$ of $\ln(1+x)$ approximates $\ln(1.1)$ with error at
        most $10^{-6}$ (true value $\ln(1.1)=0.095310$).
  \item Suppose $|f'''(t)|\leq M$ for all $t$ between $a$ and $x$.
        Derive a bound for $|f(x)-T_{2}(x)|$ in terms of $M$ and
        $|x-a|$.
\end{problems}

\section{Additional Solutions}

\begin{solutions}
  \item $f(0)=1$; $f'(x)=5(1+x)^{4}$, $f'(0)=5$. Hence
        \[
          T_{1}(x)=1+5x .
        \]

  \item $f(0)=0$; $f'(x)=\cos(x)$, $f'(0)=1$. Hence $T_{1}(x)=x$ (the
        small-angle approximation) and $\sin(0.1)\approx 0.1$. True
        value $0.099833$; the error is about $1.7\cdot 10^{-4}$
        (matching the size of the first omitted term,
        $\tfrac{x^{3}}{3!}=\tfrac{0.001}{6}$).

  \item $f(1)=1$; $f'(x)=4x^{3}$, $f'(1)=4$; $f''(x)=12x^{2}$,
        $f''(1)=12$, so $\tfrac12 f''(1)=6$. Hence
        \[
          T_{2}(x)=1+4(x-1)+6(x-1)^{2} .
        \]

  \item $f(2)=\tfrac12$; $f'(x)=-1/x^{2}$, $f'(2)=-\tfrac14$;
        $f''(x)=2/x^{3}$, $f''(2)=\tfrac28=\tfrac14$, so
        $\tfrac12 f''(2)=\tfrac18$. Hence
        \[
          T_{2}(x)=\frac12-\frac14(x-2)+\frac18(x-2)^{2} .
        \]

  \item All derivatives equal $1$ at $0$, so
        \[
          T_{3}(x)=1+x+\frac{x^{2}}{2}+\frac{x^{3}}{6} .
        \]
        At $x=0.2$:
        \[
          T_{3}(0.2)=1+0.2+\frac{0.04}{2}+\frac{0.008}{6}
                    =1+0.2+0.02+0.001333=1.221333 .
        \]
        True value $1.221403$; the error is about $6.9\cdot 10^{-5}$.

  \item $f(0)=0$; $f'(x)=\dfrac{2}{1+2x}$, $f'(0)=2$;
        $f''(x)=-\dfrac{4}{(1+2x)^{2}}$, $f''(0)=-4$, so
        $\tfrac12 f''(0)=-2$. Hence
        \[
          T_{2}(x)=2x-2x^{2} .
        \]
        (Equivalently: substitute $u=2x$ into
        $\ln(1+u)\approx u-\tfrac12 u^{2}$.)

  \item $f(1)=1$; $f'(K)=0.3\,K^{-0.7}$, $f'(1)=0.3$;
        $f''(K)=0.3\cdot(-0.7)\,K^{-1.7}=-0.21\,K^{-1.7}$,
        $f''(1)=-0.21$, so $\tfrac12 f''(1)=-0.105$. Hence
        \[
          T_{2}(K)=1+0.3\,(K-1)-0.105\,(K-1)^{2} .
        \]
        At $K=1.1$:
        \[
          T_{2}(1.1)=1+0.03-0.105\cdot 0.01=1.03-0.00105=1.02895 .
        \]
        True value $1.029006$; the error is about $5.6\cdot 10^{-5}$.

  \item Here $n=1$, $a=1$, $x=1.05$ and $f''(t)=-1/t^{2}$. On
        $[1,1.05]$, $|f''(t)|=1/t^{2}\leq 1=:M$. The remainder bound
        gives
        \[
          |R_{1}(1.05)|\leq\frac{M}{2!}\,(0.05)^{2}
            =\frac{0.0025}{2}=1.25\cdot 10^{-3} .
        \]
        Actual error: $T_{1}(1.05)=0.05$ versus $0.048790$, i.e.\
        $1.21\cdot 10^{-3}$, close to the bound, because
        $1/t^{2}\approx 1$ throughout the interval.

  \item For $f(x)=\ln(1+x)$,
        $f^{(n+1)}(t)=\dfrac{(-1)^{n}\,n!}{(1+t)^{n+1}}$, so on
        $[0,0.1]$ we have $|f^{(n+1)}(t)|\leq n!$ and
        \[
          |R_{n}(0.1)|\leq\frac{n!}{(n+1)!}\,(0.1)^{n+1}
            =\frac{(0.1)^{n+1}}{n+1} .
        \]
        Tabulate:
        \[
          n=4:\ \frac{(0.1)^{5}}{5}=2\cdot 10^{-6},
          \qquad
          n=5:\ \frac{(0.1)^{6}}{6}\approx 1.7\cdot 10^{-7} .
        \]
        The first value below $10^{-6}$ occurs at $n=5$. Check against
        the true value $0.095310$: $T_{4}(0.1)=0.1-0.005+0.000333
        -0.000025=0.095308$ (error $1.9\cdot 10^{-6}$, indeed too
        large), while $T_{5}(0.1)=0.095308+0.000002=0.095310$ (error
        about $1.5\cdot 10^{-7}$). So $n=5$.

  \item Apply the Lagrange form with $n=2$: there is $\xi$ between $a$
        and $x$ with
        \[
          f(x)-T_{2}(x)=\frac{f'''(\xi)}{3!}(x-a)^{3},
        \]
        hence
        \[
          |f(x)-T_{2}(x)|\leq\frac{|f'''(\xi)|}{6}\,|x-a|^{3}
            \leq\frac{M}{6}\,|x-a|^{3} .
        \]
\end{solutions}
